

% ---------------------------------------------------------------
%  SLIDE 1 — Overview of the controversy
% ---------------------------------------------------------------
\begin{frame}{The Controversy: To Transform or Not?}

\begin{columns}[T]
  \begin{column}{0.48\textwidth}
    \begin{block}{The CoDa School (Aitchison)}
      \begin{itemize}
        \item Work \textbf{exclusively} on log-ratio transformed data
        \item Imposes strict \textbf{Principles}: scale invariance,
              subcompositional coherence, permutation invariance
        \item Any method not satisfying these Principles is
              considered \textbf{invalid}
      \end{itemize}
    \end{block}

    \vspace{0.3cm}
    \begin{block}{The Alternative View}
      \begin{itemize}
        \item Work \textbf{directly on the simplex}, or use other
              transformations ($\alpha$-transformation, square-root, \ldots)
        \item The Principles are \textbf{neither necessary nor
              always appropriate}
        \item Choice of method should be \textbf{data-driven} and
              problem-dependent
      \end{itemize}
    \end{block}
  \end{column}

  \begin{column}{0.48\textwidth}
    \begin{alertblock}{Why Does This Matter?}
      The CoDa Principles act as a \textbf{gatekeeping criterion}:
      valid and useful methods have been dismissed simply for not
      conforming to them.

      \vspace{0.2cm}
      This is \textbf{unusual, perhaps unique, in statistical
      data analysis} (Scealy \& Welsh 2014).
    \end{alertblock}

    \vspace{0.3cm}
    \begin{block}{Key Questions}
      \begin{enumerate}
        \item Is \textbf{scale invariance} always meaningful?
        \item Is \textbf{subcompositional coherence} always required?
        \item Should the Principles \textbf{exclude} alternative methods?
      \end{enumerate}
    \end{block}
  \end{column}
\end{columns}

\end{frame}

% ---------------------------------------------------------------
%  SLIDE 2 — Watson
% ---------------------------------------------------------------
\begin{frame}{Key Critics: D.F.\ Watson}

\begin{columns}[T]
  \begin{column}{0.55\textwidth}
    \textbf{Watson \& Philip (1989)} were among the \textbf{earliest} voices
    to challenge the log-ratio framework in the pages of
    \textit{Mathematical Geology}.

    \vspace{0.3cm}
    \begin{block}{Core Argument}
      \begin{itemize}
        \item Questioned whether \textbf{scale invariance} is a
              fundamental requirement for compositional analysis
        \item If compositions are strictly defined as summing to one,
              scale invariance is \textbf{already built in} — and imposing
              it as a separate principle becomes circular
        \item Their exchanges with Aitchison (1989--1992) were
              instrumental in \emph{sharpening} the CoDa Principles,
              even if from an adversarial position
      \end{itemize}
    \end{block}
  \end{column}

  \begin{column}{0.42\textwidth}
    \begin{alertblock}{Historical Note}
      The debate was cut short by a notable incident: Aitchison
      used \textbf{derogatory language} toward Watson in the
      published exchange, which effectively ended Watson's
      participation.

      \vspace{0.2cm}
      This incident itself illustrates the \textbf{heated and
      personal} nature of the controversy.
    \end{alertblock}

    \vspace{0.3cm}
    \begin{block}{Legacy}
      Watson's questions about scale invariance remain
      \textbf{unresolved} and were later taken up more formally
      by Scealy \& Welsh (2011, 2014).
    \end{block}
  \end{column}
\end{columns}

\end{frame}

% ---------------------------------------------------------------
%  SLIDE 3 — Zier & Rehder
% ---------------------------------------------------------------
\begin{frame}{Key Critics: Zier \& Rehder (1998, 2001)}

\begin{columns}[T]
  \begin{column}{0.55\textwidth}
    Zier \& Rehder continued the challenge against the CoDa School,
    focusing on \textbf{technical and foundational} issues.

    \vspace{0.3cm}
    \begin{block}{Core Arguments}
      \begin{itemize}
        \item Challenged the mathematical definition of a
              \textbf{metric} within compositional data analysis
        \item Argued that scale invariance is \textbf{problematic}
              when compositions are strictly constrained to sum to one
        \item Pointed out that even \textbf{Aitchison's own earlier
              work} did not strictly adhere to the current CoDa
              definition of scale invariance
        \item Made the case that log-ratio methods are not the
              \emph{only} valid approach
      \end{itemize}
    \end{block}
  \end{column}

  \begin{column}{0.42\textwidth}
    \begin{block}{Reception by the CoDa School}
      Despite raising \textbf{some valid points}, their criticisms
      were \textit{dismissed relatively easily} — partly because
      of technical errors in their own work that the CoDa School
      used to deflect the substantive arguments.
    \end{block}

    \vspace{0.3cm}
    \begin{alertblock}{Broader Pattern}
      This illustrates a recurring dynamic: \textbf{genuine
      methodological concerns} are sidelined when the critic
      makes any technical mistake, however minor relative to
      the central point.
    \end{alertblock}
  \end{column}
\end{columns}

\end{frame}

% ---------------------------------------------------------------
%  SLIDE 4 — Butler & Glasbey
% ---------------------------------------------------------------
\begin{frame}{Key Critics: Butler \& Glasbey (2008)}

\begin{columns}[T]
  \begin{column}{0.55\textwidth}
    Butler \& Glasbey represent a different kind of challenge:
    \textbf{practical, application-driven} work that did not
    conform to the CoDa Principles.

    \vspace{0.3cm}
    \begin{block}{Their Approach}
      \begin{itemize}
        \item Developed a \textbf{latent Gaussian model} — an
              otherwise sensible method — that projects a
              multivariate Gaussian onto the unit simplex
        \item The method handles \textbf{zero values} naturally,
              a genuine limitation of the log-ratio framework
        \item Useful, interpretable, and practically motivated
      \end{itemize}
    \end{block}

    \vspace{0.2cm}
    \begin{alertblock}{The CoDa School's Response}
      They were required to \textbf{list and explain their
      ``failings''} relative to the Principles, and their method
      was downgraded to a \emph{diagnostic or exploratory tool}
      — not a full analysis method.
    \end{alertblock}
  \end{column}

  \begin{column}{0.42\textwidth}
    \begin{block}{What This Reveals}
      \begin{itemize}
        \item The Principles function as a \textbf{barrier to
              methodological innovation}
        \item A method can be \textbf{statistically sound and
              practically useful} yet be excluded from the CoDa
              canon
        \item The question is not \emph{does it work?} but
              \emph{does it comply?}
      \end{itemize}
    \end{block}

    \vspace{0.3cm}
    \begin{block}{Zero Problem}
      Ironically, the log-ratio approach cannot handle zeros at
      all, while Butler \& Glasbey's ``non-compliant'' method
      does — a \textbf{practical advantage} dismissed on
      theoretical grounds.
    \end{block}
  \end{column}
\end{columns}

\end{frame}

% ---------------------------------------------------------------
%  SLIDE 5 — Stewart, Tsagris, Scealy & Welsh
% ---------------------------------------------------------------
\begin{frame}{Key Critics: Stewart, Tsagris, Scealy \& Welsh}

\begin{columns}[T]
  \begin{column}{0.55\textwidth}
    A cluster of influential papers from \textbf{2010--2014}
    provided both empirical and theoretical challenges.

    \vspace{0.3cm}
    \begin{block}{Stewart \& Field (2010); Tsagris, Preston \& Wood (2011)}
      \begin{itemize}
        \item Provided \textbf{real data examples} showing that
              different methods can yield \textbf{substantively
              different conclusions}
        \item Differences are largest for \textbf{highly variable
              data near or on the boundary of the simplex} — exactly
              where log-ratio methods struggle most (zeros,
              near-zeros)
        \item Proposed the \textbf{$\alpha$-transformation} as a
              flexible, data-driven alternative
      \end{itemize}
    \end{block}
  \end{column}

  \begin{column}{0.42\textwidth}
    \begin{block}{Scealy \& Welsh (2011a, 2014)}
      The most \textbf{theoretically rigorous} critique:
      \begin{itemize}
        \item Challenged the \textbf{logical motivation} for scale
              invariance and subcompositional coherence
        \item Proposed \textbf{amalgamations} as more natural
              marginals for compositional data
        \item Applied the \textbf{Kent distribution} on the
              hypersphere — a principled alternative geometry
        \item Argued that the Principles are \textbf{not logically
              necessary} for valid inference
      \end{itemize}
    \end{block}

    \vspace{0.1cm}
    \begin{alertblock}{}
      Their work shifts the debate from \emph{practice} to
      \textbf{foundations}.
    \end{alertblock}
  \end{column}
\end{columns}

\end{frame}

% ---------------------------------------------------------------
%  SLIDE 6 — Conclusion on the controversy
% ---------------------------------------------------------------
\begin{frame}{Where Does the Controversy Stand?}

\begin{columns}[T]
  \begin{column}{0.48\textwidth}
    \begin{block}{What the Critics Agree On}
      \begin{enumerate}
        \item The definition of a \textbf{composition} is contested
        \item \textbf{Scale invariance} is not always meaningful or
              necessary
        \item \textbf{Subcompositional coherence} has questionable
              logical foundations
        \item The CoDa Principles \textbf{actively exclude} valid
              alternative methods
        \item The log-ratio approach \textbf{fails with zeros} and
              near-boundary data
      \end{enumerate}
    \end{block}
  \end{column}

  \begin{column}{0.48\textwidth}
    \begin{block}{The Current Landscape}
      \begin{itemize}
        \item No consensus — the debate \textbf{remains open}
        \item A growing body of work (Tsagris, Scealy, Stewart,
              \ldots) offers \textbf{principled alternatives}
        \item Practitioners increasingly \textbf{try multiple
              methods} and let the data decide
      \end{itemize}
    \end{block}

    \vspace{0.3cm}
    \begin{alertblock}{Our Position}
      We take the view that the \textbf{choice of method should
      be data-driven}. We will explore the $\alpha$-transformation
      and random forest models as flexible tools that do not
      require adherence to the CoDa Principles.
    \end{alertblock}
  \end{column}
\end{columns}

\end{frame}

